\section{本章小结}
\label{sec:heat-summary}

本章回答了一个局部规律怎样决定整体演化的问题。每一点向邻居的平均值靠拢，写成方程是$\partial_tu=\Delta u+f$，再加上一个初值和边界条件，问题就确定了。固定温度和无通量是两种不同的物理规定，前者让热量流出，后者保留一个零模式。

同一个传播过程有四种描述，它们各有适用的场合，见表~\ref{tab:heat-summary}。边界模式把问题对角化，高频衰减得快，零模式不衰减；热核把一点的影响写成Gauss钟形，宽度按$\sqrt{t}$增长；能量和极值估计不需要显式公式；半群把解看成算子，边界条件落在生成元的定义域中。

\begin{BookInlineTable}
  \centering
  \caption{同一扩散过程的四种描述。}
  \label{tab:heat-summary}
  \begin{tabular}{lll}
    \toprule
    描述 & 适用情形 & 看到的结构\\
    \midrule
    边界模式 & 有界区域，零边值 & 独立模式按$\mathrm{e}^{-\lambda_kt}$衰减\\
    热核卷积 & 全空间，受增长限制 & Gauss平均，$\sqrt{t}$尺度\\
    能量与极值 & 一般区域，明确的边界 & 耗散、比较、唯一性\\
    解半群 & 选定函数空间 & 时间拼接，无界生成元\\
    \bottomrule
  \end{tabular}
\end{BookInlineTable}

结论的成立范围都有明确的边界。初值恢复的方式取决于数据：有界一致连续初值一致恢复，$L^p$初值（$1\leq p<\infty$）在$L^p$中恢复，一般$L^\infty$初值只在连续点逐点恢复。全空间的唯一性要限制增长，否则有Tikhonov解。经典解的公式需要相应的正则性和相容性，粗糙数据时只能在$L^2$或弱的意义下理解。估计都是有限时间的，长期行为的一般理论不在本卷。

持续输入和非齐次边界都归结到同一条原则：线性叠加，把输入切片后逐片自由演化。这个原则要求已经知道自由演化；非零边值要先提升，全空间卷积不能直接用于有界区域。

平滑效应是扩散最突出的特点，也是逆问题困难的根源。正向过程在$t>0$时把高频按指数压低，反向恢复必须把它们放大回来，所以可恢复的频率只随测量精度的对数增长。

第~\ref{chap:wave-equations}章的波动方程有两个初值、有限的传播速度和守恒的能量，可以与这里的耗散和瞬间影响逐条比较。随机路径、转移概率和Fokker--Planck方程给出的概率解释不属于本卷范围，空间离散则在第~\ref{chap:pde-numerical-methods}章讨论。
